\addchap{Introduction} 



%\addsec{Chapter Introduction (84 words)}
Uncrewed systems have made exposure one of the key problems in contemporary land warfare. The literature shows that uncrewed systems are reshaping reconnaissance, targeting and strike, while wider scholarship on military innovation and mission command explains adaptation under changing conditions. However, these literatures rarely connect. Hence the implications for light-infantry tactics, small-unit command and force design remain underexplored. This thesis addresses that lacuna through qualitative synthesis of scholarly and doctrinal sources, practitioner interviews and systematic analysis of open-source combat evidence from the Russo--Ukrainian War. It identifies recurring patterns in order to derive implications for light infantry adaptation.

%\addsec{Literature Review (873 words)}
Ukraine-focused literature establishes that \gls{uxs} accelerate \gls{killchain}, extend \gls{isr} and tighten fires integration, but leaves unresolved how those effects govern light-infantry movement, mission command and small-state force design \parencite{MOLLOY_2024,WATLING_2025B,BRONK_2025,CORNISH_2024}.


%Recent secondary literature based on Ukraine-focused fieldwork and analysis by \gls{rusi} researchers and \textcite{MOLLOY_2024} provides a detailed account of how \gls{uxs} are reshaping the character of warfare. These studies highlight accelerated kill-chains, persistent ISR and the tighter integration of fires as defining features of the contemporary battlefield. However, the very factors which make this literature valuable --- timeliness, privileged access and operational security constraints --- also reduce transparency and limit independent auditability. As a result, while the literature offers rich descriptive insight, its conclusions are often difficult to generalise beyond the Russo--Ukrainian War.

%At the same time, the empirical record complicates technologically deterministic interpretations. It is tempting to conclude that UAS will displace artillery and combined-arms manoeuvre. Yet recent studies consistently emphasise the continued primacy of indirect fires \parencite{BRONK_2022a,BRONK_2025,CORNISH_2024,MOLLOY_2024,REYNOLDS_2023B,WATLING_2025B}. \textcite{vers_2024} demonstrates that even platoon-level manoeuvre depends on the tightly coordinated integration of EW, fires, ISR, engineers and drones, increasing the complexity of offensive action rather than simplifying it. This is consistent with a broader literature on military innovation, which emphasises the decisive role of organisational adaptation over technological invention alone \parencite{KREPINEVICH_1994,OCHS_2020,OWENS_2002}. Furthermore, the Russo--Ukrainian War has largely become positional rather than manoeuvre-centric, adding further nuance to apparent drone-driven change.

%Within this framework, UAS are best understood as enablers of existing fires systems rather than replacements for them. \gls{term:isrd}s tighten the \gls{killchain}, providing persistent targeting and immediate damage assessment \parencite{BRONK_2024,BRONK_2025,WATLING_2025B}. While precision rotary-wing \gls{term:owad} payloads are often insufficient to destroy armoured platforms outright and may suffer high attrition rates \parencite{WATLING_2025B}, their primary effect lies in immobilisation and \textit{fixing}. Fixed- and rotary-wing systems frequently enable artillery to complete the kill, while also proving highly effective against light infantry and combat service support. However, this raises unresolved questions regarding lethality attribution and battlefield effectiveness. As \textcite{MOLLOY_2024} notes, apparent drone effectiveness reflects compensation for capability shortfalls or attribution bias in video-heavy reporting. Even where UAS are assessed to account for a significant proportion of losses \parencite{WATLING_2025B}, this has not translated into decisive battlefield outcomes. The literature does not resolve these ambiguities, nor does it examine their implications for the infantry's evolving role within the kill-chain.

At the tactical level, the migration of reconnaissance to drones \parencite{WATLING_2023A} has altered the conditions of survivability. Infantry can no longer rely on concealment in proximity to the \gls{flot}. Instead, survivability increasingly depends on dispersion, fortification and \gls{ems} discipline \parencite{WATLING_2023A,WATLING_2025B,WATLING_2023B}. Since UAS extend the EMS into the frontline, disciplined signal emission and effective \gls{ew} are required, with \gls{elelcrecce} by overflying systems becoming more common \parencite{WATLING_2025A}. These conditions sharpen the trade-off between communication and detection, placing greater reliance on decentralised execution at lower levels.

However, pervasive ISR simultaneously generates a countervailing tendency. As \textcite{COHEN_1996} observed, new technologies enable senior commanders to ``perch cybernetically'' alongside subordinates. Persistent surveillance intensifies this dynamic, increasing the temptation to centralise control even as dispersion demands greater autonomy. The literature recognises both trends but does not adequately reconcile them, leaving it unclear how mission command functions in practice under conditions of persistent observation and electromagnetic exposure.

Conceptually, early work anticipated some of these dynamics. \textcite[10]{Noel_2013_10} suggested that drones could \textit{occupy} the land from the air. Contemporary evidence indicates that their effect is better understood as a form of aerial denial \parencite{Stepanenko_2025b}, blurring the concept of the FLOT and constraining manoeuvre and sustainment. Yet the implications of this shift for small-unit tactics remain under-theorised.

The literature on \gls{ugs} further complicates the picture. Rather than mirroring aerial systems, \textcite{HINTON_2023} argues that their near-term value lies in exposure-reducing support roles, given constraints in mobility, power and teleoperation. As a result, UGS remain comparatively under-explored in both their tactical impact and their implications for light-infantry organisation.

A longstanding sceptical tradition cautions against overstating the transformative impact of new technologies. \textcite{KREPINEVICH_1994} argues that invention without organisational adaptation is ineffective, while \textcite{BETTS_1996} highlights the tendency of militaries to misuse new capabilities when doctrine and culture lag behind. More broadly, critics of the Revolution in Military Affairs emphasise evolutionary adaptation over rupture \parencite[50--52]{ALACH_2008}, with drone warfare itself often characterised as incremental rather than transformative \parencite{RASSLER_2015}. \textcite{kunertova_2024_2024_08} extends this critique directly to Ukraine, arguing that observed innovation reflects improvisation under constraint rather than a technological revolution, while \textcite{hinote_ryan_2024} emphasise that drones are only transformative when integrated within effective C2 and sensor networks.

%Taken together, the literature provides a detailed but fragmented picture. It establishes what UxS enable tactically and highlights important trends in survivability, fires integration and command. However, it leaves unresolved how these dynamics translate into small-unit behaviour, how they reshape the practice of mission command under persistent ISR, and how tactical effects scale into organisational and strategic outcomes. These gaps form the point of departure for the present study.


%\addsec{Research Lacunae: So what for this thesis? (353 words)}

%The literature establishes what UAS enable tactically --- accelerated kill-chains, persistent ISR and tighter integration of fires. It leaves three critical gaps for light infantry adaptation.

The generalisability of Russo--Ukrainian observations remains contested. Western doctrine, force protection measures and EW capabilities may mitigate some threats, but this assumption has not been tested --- particularly at the sub-company level. That leaves the impact of pervasive ISR on section and platoon behaviour unresolved. Ukraine-focused fieldwork documents a shift toward dispersion, fortification and electromagnetic discipline. It does not derive how these principles translate into section \gls{battledrill}. For example, while the ISR mechanism is well documented, its translation into Western doctrine remains underdeveloped.

Pervasive ISR also creates a tension between decentralised execution and the centralising pull of ubiquitous observation. It allows senior commanders to interrogate subordinate action while dispersed units require greater autonomy. Western doctrine emphasises mission command, but is silent on how junior leader initiative operates when multiple echelons observe the same fight. The literature therefore leaves mission command under persistent surveillance insufficiently explained.

Existing studies also treat battlefield effectiveness, organisational adaptation and strategic implications separately. They do not explain what drone warfare means for resource-constrained forces. In particular, they do not show how UxS alter \gls{rcp} ratios. This gap is especially acute for small states such as Ireland, whose strategic culture, resource constraints and force structure differ markedly from those of larger powers. Given the operational sensitivity of TTP evolution, this thesis adopts an indirect analytical approach, examining observable patterns, organisational responses and doctrinal implications at small-unit level through three questions:



\begin{enumerate}
	\item How have uncrewed systems reshaped the contemporary battlefield, particularly in terms of survivability, sensing, movement and the physical conditions of light infantry tactics?
	\item How do persistent ISR and drone-enabled lethality alter light infantry tactics, mission command, platoon composition and tactical adaptation at sub-company level?
	\item What do these tactical and command changes imply for Irish force design, military utility, small-state RCP and deterrence-by-denial\index{Deterrence-by-Denial}?
\end{enumerate}


%\addsec{Sources (365 words)}
%Primary sources form the empirical foundation. Doctrine, training publications and capability statements from relevant militaries are used selectively to contextualise how mission command, TTPs and UAS/UGS employment are framed. Open-source material from the Russo--Ukrainian War --- including official communiqu\'{e}s, imagery, combat footage and interviews --- provides the core empirical base. Approximately 2,000 primary-source items were reviewed, predominantly combat footage from verified institutional Telegram channels and verified combatant sources. From this corpus, 178 sources were selected, verified and subjected to systematic bilateral thematic analysis, examining Ukrainian and Russian material in parallel and coding for recurring behavioural patterns across survivability, movement, fires coordination, sustainment and command. This is not a sample of convenience. It represents sustained analytical engagement with the primary record of this war at section and platoon level. These sources are treated as partial and success-biased; the analytical question is not whether a unit published all available footage, but whether what was published authentically documents the tactical phenomenon under examination.

%Secondary sources provide the main analytical engine. Peer-reviewed scholarship supplies conceptual frameworks for assessing whether drones drive genuine tactical change or amplify an enduring character. Practitioner-oriented analysis offers mechanism-rich accounts of ISR saturation, EW friction and the evolving hider--finder contest. Irish and EU policy documents provide the contextual constraints for a small, professional army.

%Semi-structured interviews with a small number of practitioners and subject-matter experts complement documentary sources. This captures how junior initiative, information discipline and organisational friction are experienced in practice and helps to contextualise theatre-specific evidence.



%\addsec{Methodology (284 words)}
The empirical foundation is a bilateral thematic analysis of approximately 2,000 primary-source items --- combat footage, communiqu\'{e}s, imagery and interviews --- from the Russo--Ukrainian War. From this corpus, 178 sources were selected, verified and analysed thematically \parencite{BRAUN_2006}, coding observable actions into recurring behavioural patterns across survivability, movement, fires coordination, sustainment and command. Examining Ukrainian and Russian material in parallel guards against single-belligerent bias and strengthens pattern validity across the corpus \parencite{WEBB_1966,JICK_1979}. These sources are treated as partial and success-biased; the analytical question is whether published footage authentically documents the tactical phenomenon under examination, not whether it is statistically representative \parencite{CASTRO_2013}. Secondary sources --- peer-reviewed scholarship, practitioner analysis and Irish policy documents --- supply conceptual frameworks. Semi-structured interviews with practitioners and subject-matter experts, analysed through an essentialist lens \parencite{FINNEGAN_2013}, contextualise theatre-specific evidence and capture how junior initiative and organisational friction operate in practice. Ethical approval and informed consent governed data collection.



%This thesis uses a qualitative methodology based on triangulating written sources, visual open-source material and interviews, consistent with established within-method approaches which strengthen inference where single-method designs are insufficient and replication is constrained \parencite{WEBB_1966,JICK_1979}. The purpose is to identify recurring patterns in how light infantry units adapt to the presence of uncrewed systems. It emphasises mechanism and convergence across sources rather than measurement or frequency --- appropriate in military contexts characterised by operational security, partial disclosure and institutional opacity \parencite{CASTRO_2013,JICK_1979}.
%Combat footage undergoes reflexive thematic analysis \textcite{BRAUN_2006}, which suits heterogeneous visual and interview data for systematically identifying adaptation patterns. Initial codes capture observable actions and are iteratively grouped into tactical themes, with credibility enhanced through source triangulation. The bilateral structure --- examining Ukrainian and Russian material in parallel --- guards against single-belligerent bias and strengthens pattern validity across the corpus.
%Interviews are conducted with a small number of practitioners, journalists and subject-matter experts and are analysed semantically through an essentialist/realist lens \parencite{FINNEGAN_2013}. They validate interpretations drawn from open sources, clarify how decisions are made and delegated in practice, and identify institutional constraints relevant to an Irish context.
%The study seeks inductive-analytical rather than statistical generalisation. Ethical approval, informed consent and operational security considerations shape data collection and reporting. Scope is delimited to contemporary Europe and the Russo--Ukrainian War. Legal and moral considerations --- such as determining human intent and surrender --- fall outside the scope. English-language, land-component secondary sources dominate, reflecting the author's positionality within the Irish Army.

%Direct involvement in UAS procurement and doctrine development created a risk of confirmation bias, mitigated in three ways: sceptical voices were deliberately engaged, including Kunertova (2024), Betts (1996) and Molloy (2024)'s caveats on attribution bias; the bilateral analytical structure required observed patterns to hold across both Ukrainian and Russian material; and Irish procurement knowledge was treated as a hypothesis to be tested rather than a conclusion to be confirmed.

The Russo--Ukrainian War is the principal case because it provides the most demanding observable environment --- high-intensity, adaptive and technologically contested. Alternative theatres were considered but offer inferior volume, quality and variety. While \textcite{watling_2025_10} cautions that its conditions may not translate universally, the underlying mechanisms nevertheless provide analytically transferable lessons. While Ukrainian and Western sources dominate, Russian footage and official commentary were analysed throughout. Open-source video material is subject to selection and success bias and cannot provide statistically representative data. However, the scale, recurrence and cross-source consistency of observed behaviours allow for the identification of tactically significant patterns. This thesis therefore treats such material not as quantitative evidence, but as a qualitative corpus from which recurring adaptations in survivability, fires integration and small-unit behaviour can be derived.

The corpus is deliberately land-focused because the thesis examines light-infantry adaptation. Its main limitation is that the Russo--Ukrainian War is ongoing and many relevant lessons are classified, or not yet codified in doctrine, capability codes or capability statements. Hence, institutional outputs lag battlefield practice. The author’s position as a serving Irish Army officer also shapes the analytical frame, particularly the emphasis on section-level adaptation, mission command and small-state force design. The thesis therefore draws conclusions from observable tactical behaviour, practitioner evidence and open-source documentation. It uses bilateral Ukrainian/Russian coding, Cyrillic primary sources and triangulation to identify recurring mechanisms. This thesis is intended to be actively disseminated among infantry units, junior commanders and Defence Forces' staff officers concerned with the implications of persistent ISR for light infantry, force design and deterrence by denial.

%\addsec{Thesis Outline (101 words)}
Chapter one shows how UxS have democratised ISR and precision strike, making exposure a core tactical problem. It establishes the battlefield conditions on which the rest of the thesis builds. Chapter two shows how persistent surveillance changes light-infantry tactics, command and organisation. It develops the concepts of movement as timed exposure, the infantry ceiling and conditional manoeuvre, before proposing adapted section structures built around organic ISR, \gls{cuas} and stronger junior leadership. Chapter three examines how drone-enabled mass and cost asymmetry reshape force design and deterrence. It argues that small states can generate disproportionate RCP through integrated UxS, enabling credible deterrence-by-denial under resource constraints.
